We investigate how transformers internally represent integer intervals, number comparisons, and modular arithmetic in their residual streams. Using linear probing, regression probes, and activation patching on GPT-2 Small, we find that the model stores rich numerical information in its hidden states even when it cannot behaviorally solve the task. Comparison information (is $A > B$?) is linearly decodable from layer~1 onward, reaching 99.1\% accuracy by layer~9, despite the model achieving only 48.6\% behavioral accuracy. Interval membership (is $X \in [A, B]$?) peaks at 83.8\% probe accuracy at layer~8---significantly above chance but 15~percentage points below comparison, consistent with the hypothesis that intervals require composing two boundary checks. Most strikingly, ring counting answers (modular arithmetic over the alphabet) are decodable at 58\% from layer~12 (versus 3.8\% chance), despite 0\% behavioral accuracy. These findings reveal a systematic dissociation between internal computation and behavioral performance: transformers build internal representations of numerical operations that are never surfaced in their outputs. Our results suggest that behavioral evaluations alone dramatically underestimate what models internally compute, and that probing is essential for understanding transformer numerical capabilities.
